Use affine coordinates $A=(0,3)$, $B=(0,0)$, $C=(3,0)$, which preserve ratios along lines. Then $D=(1,0)$ because $BD:DC=1:2$. Since $AF:FD=1:2$, $F=(\dfrac{1}{3},2)$. Line $BF$ has equation $y=6x$, while $AC$ has equation $y=3-x$. Thus $E=(\dfrac{3}{7},\dfrac{18}{7})$, one-seventh of the way from $A$ to $C$. Hence $\boxed{AE:EC=1:6}$.
\begin{asy}
texpreamble("\usepackage{amsmath,amssymb}");
size(220);
defaultpen(fontsize(9pt));
pair A=(0,3),B=(0,0),C=(3,0); draw(A--B--C--cycle); label("$A$",A,(-1,-1)); label("$B$",B,(1,-1)); label("$C$",C,(0,1)); pair D=(1,0),F=(1/3,2),E=(3/7,18/7);draw(A--D);draw(B--E);dot(F);label("$D$",D,(0,-1));label("$F$",F,(-1,0));label("$E$",E,(1,1));
\end{asy}
